\section{Method}
\label{sec:method}

We built \sys to solve the above challenges. Given the model, workload characteristic, latency requirements, and SLO attainment target, \sys will determine (a) the parallelism strategies for prefill and decoding instances, (b) the number of each instance type to deploy, as well as (c) how to place them onto the physical cluster. We call the solution a \textit{placement}. Our goal is to find a placement that maximizes the per-gpu goodput.

As explained in \S\ref{subsec:practical_problems}, a key design consideration is to manage communications between disaggregated prefill and decoding phases, given varying cluster setups.
In this section, we first present two placement algorithms: one for clusters with high-speed cross-node networks (\S\ref{subsec:optimal_placement}) and the other for environments lacking such infrastructure (\S\ref{subsec:practical_placement}); the latter introduces additional constraints. We then develop online scheduling optimizations that adapt to the nuances of real-world workloads (\S\ref{subsec:online_scheduling}).

\subsection{Placement for High Node-Affinity Cluster}
\label{subsec:optimal_placement}

\begin{algorithm}[t!]
\caption{High Node-Affinity Placement Algorithm}\label{alg:optimal-placement}
\begin{algorithmic}
\Require LLM $G$, \#node limit per-instance $N$, \#GPU per-node $M$, GPU memory capacity $C$, workload $W$, traffic rate $R$.
\Ensure the placement $\mathit{best\_plm}.$ 
    \State $\mathit{config_p}, \mathit{config_d} \leftarrow \emptyset, \emptyset$
    \For{$\mathit{intra\_op} \in \{1, 2, ..., M\}$}
        \For{$\mathit{inter\_op} \in \{1, 2, ..., \frac{N\times M}{\mathit{intra\_op}}\}$}
            \If{$\frac{G.size}{\mathit{inter\_op} \times \mathit{intra\_op}} < C$}
                \State $\mathit{config} \leftarrow (\mathit{inter\_op}, \mathit{intra\_op})$
                \State $\hat{G} \leftarrow \text{parallel}(G, \mathit{config})$
                \State $\mathit{config.goodput} \leftarrow \text{simu\_prefill}(\hat{G}, W)$
                \If{$\frac{\mathit{config_p.goodput}}{config_p.num\_gpus} < \frac{\mathit{config.goodput}}{config.num\_gpus}$}
                    \State $\mathit{config_p} \leftarrow \mathit{config}$
                \EndIf
                \State $\mathit{config.goodput} \leftarrow \text{simu\_decode}(\hat{G}, W)$
                \If{$\frac{\mathit{config_d.goodput}}{config_d.num\_gpus} < \frac{\mathit{config.goodput}}{config.num\_gpus}$}
                    \State $\mathit{config_d} \leftarrow \mathit{config}$
                \EndIf
            \EndIf
        \EndFor
    \EndFor
    \State $n, m \leftarrow \lceil \frac{R}{\mathit{config_p.goodput}} \rceil, \lceil \frac{R}{\mathit{config_d.goodput}} \rceil$ 
    \State $\mathit{best\_plm} \leftarrow (n, \mathit{config_p}, m, \mathit{config_d})$
\State \textbf{return} $\mathit{best\_plm}$
\end{algorithmic}
\end{algorithm}

On high node-affinity clusters equipped with Infiniband, KV caches transmission overhead across nodes is negligible, \sys can deploy prefill and decoding instances across any two nodes without constraints. 
We propose a two-level placement algorithm for such scenarios:  we first optimize the parallelism configurations for prefill and decoding instances separately to attain phase-level optimal per-gpu goodput; then, we use replication to match the overall traffic rate.

However, finding the optimal parallel configuration for a single instance type, such as for the prefill instance, is still challenging, due to the lack of a simple analytical formula to calculate the SLO attainment (a.k.a., percentage of requests that meet TTFT requirement), given that the workload has diverse input, output lengths, and irregular arrival patterns. Gauging the SLO via real-testbed profiling is time-prohibitive. We thus resort to building a simulator to estimate the SLO attainment, assuming prior knowledge of the workload's arrival process and input and output length distributions.
Although short-term interval is impossible to predict, the workload pattern over longer timescales (e.g.,
hours or days) is often predictable~\cite{li2023alpaserve, zhangshepherd}. \sys fits a distribution from the history request traces and resamples new traces from the distribution as the input workload to the simulator to compute the SLO attainment. Next, \sys simply enumerates the placements and finds the maximum rate that meets the SLO attainment target with binary search and simulation trials.


Algorithm~\ref{alg:optimal-placement} outlines the process. We enumerate all feasible parallel configurations, subject to cluster capacity limit, for both prefill and decoding instances. Then, for a specific prefill phase configuration, we use \verb|simu_prefill| to simulate and find its maximum goodput via binary search (similarly for using \verb|simu_decode|  for decoding). 
After determining the optimal parallel configurations for both prefill and decoding instances, we replicate them to achieve the user-required overall traffic rate according to their goodput. 

The complexity of Algorithm~\ref{alg:optimal-placement} is $O(NM^2)$, with $N$ as the node limit per instance and $M$ representing the typical number of GPUs per node in modern clusters (e.g., 8). The search space is manageable and the solving time is under 1.3 minutes in our largest setting, as demonstrated in \S\ref{exp:run_time}.

\parabf{Simulator building.} Algorithm~\ref{alg:optimal-placement} relies on a simulator to estimate the goodput under various SLOs and SLO attainment goals given the workload and the parallelism plan. 
To build an accurate simulator, we analyze the FLOPs and the number of memory accesses for prefill and decoding phases respectively, and use a latency model to approximate the inference execution time. See details in Appendix~\ref{appendix:latency_model}. The simulator aligns well with real profiling results, thanks to the high predictability of DNN workloads~\cite{clockwork, li2023alpaserve}, verified in \S\ref{exp:ablation}. 

By far, we have developed Algorithm~\ref{alg:optimal-placement} assuming we can place the prefill and decoding instance between any two nodes (or on the same node) of the cluster, and the KV cache transmission utilizes high bandwidth network. In many real clusters, GPUs inside a node access to high-bandwidth NVLINK while GPUs distributed across nodes have limited bandwidth. We next develop an algorithm to address this constraint.

\begin{algorithm}[t!]
\caption{Low Node-Affinity Placement Algorithm}\label{alg:low-affinity-placement}
\begin{algorithmic}
\Require LLM $G$, \#node limit per-instance $N$, \#GPU per-node $M$, GPU memory capacity $C$, workload $W$, traffic rate $R$.
\Ensure the placement $\mathit{best\_plm}.$ 
    \State $\mathit{config}^{*} \leftarrow \emptyset$
    \For{$\mathit{inter\_op} \in \{1, 2, ..., N\}$}
        \State $\mathcal{P} \leftarrow \text{get\_intra\_node\_configs}(G, M, C, \mathit{inter\_op})$
        \For{$P_p \in \mathcal{P}$}
            % \State $\hat{G}_p \leftarrow \text{parallel}(G, P_p, \mathit{inter\_op})$
            \For{$P_d \in \mathcal{P}$}
                \If{$P_p.\mathit{num\_gpus} + P_d.\mathit{num\_gpus} \le M$}
                    \State $\mathit{config} \leftarrow (\mathit{inter\_op}, P_p, P_d)$
                    \State $\hat{G}_p, \hat{G}_d \leftarrow \text{parallel}(G, \mathit{config})$
                    \State $\mathit{config.goodput} \leftarrow \text{simulate}(\hat{G}_p, \hat{G}_d, W)$
                        \If{$\frac{\mathit{config.^{*}goodput}}{\mathit{config.^{*}num\_gpus}} < \frac{\mathit{config.goodput}}{\mathit{config.num\_gpus}}$}
                            \State $\mathit{config^{*}} \leftarrow \mathit{config}$
                        \EndIf
                \EndIf
            \EndFor
        \EndFor
    \EndFor
    
    \State $n \leftarrow \lceil \frac{R}{\mathit{config.^{*}goodput}} \rceil$ 
    \State $\mathit{best\_plm} \leftarrow (n, \mathit{config^{*}})$

\State \textbf{return} $\mathit{best\_plm}$
\end{algorithmic}
\end{algorithm}

\subsection{Placement for Low Node-Affinity Cluster}
\label{subsec:practical_placement}
A straightforward solution is to always colocate prefill and decoding instances on the same node, utilizing the NVLINK, which is commonly available inside a GPU node.
For large models, e.g. with 175B parameters (350GB), we may be unable to even host a single pair of prefill and decoding instances in an 8-GPU node ($80G \times 8 = 640G < 350 \times 2 GB $). We incorporate this as additional placement constraints and co-optimize it with model parallelism, presented in Algorithm~\ref{alg:low-affinity-placement}.


The key insight is that KV cache transfer occurs exclusively between corresponding layers of prefill and decoding instances.
Leveraging inter-op parallelism, we group layers into stages and divide each instance into segments, termed as \textit{instance segments}, with each segment maintaining one specific inter-op stage. 
By colocating prefill and decoding segments of the same stage within a single node, we force the transfer of intermediate states to occur only via NVLINK. Inside a node, we set the same parallelism and resource allocation for segments of the same instance. Given the typical limitation of GPUs per node (usually 8), we can enumerate possible configurations inside one node and use the simulator to identify the configurations that yield the best goodput.

As outlined in Algorithm~\ref{alg:low-affinity-placement}, we begin by enumerating inter-op parallelism degrees to get all the possible instance segments. For each segment, we get all possible intra-node parallelism configurations by calling \verb|get_intra_node_configs|. Then we use simulation to find the optimal one and replicate it to satisfy the target traffic rate.

\subsection{Online scheduling}
\label{subsec:online_scheduling}


The runtime architecture of \sys is shown in Figure~\ref{fig:architecture}. \sys operates with a simple FCFS scheduling policy. All incoming requests arrive at a centralized controller, then dispatched to the prefill instance with the shortest queue for prefill processing, followed by dispatch to the least loaded decoding instance for decoding steps. This setup, while simple, is optimized with several key enhancements tailored to the nuances of real-world workloads.

\parabf{Reducing pipeline bubbles.} 
To mitigate the pipeline bubbles caused by non-uniform prompt lengths (\S\ref{subsec:practical_problems}), we schedule the requests in a way that balances the execution time across all batches in the pipeline. This is achieved by noting that, for both prefill and decoding instances, the number of new tokens in the batch is a reliable indicator of the batch's real execution time.
For prefill instances, we profile the target model and GPU to figure out the shortest prompt length $L_m$ needed to saturate the GPU. We schedule prefill batches with a total sequence length close to $L_m$, by either batching multiple requests shorter than $L_m$ or individually scheduling requests longer than $L_m$. For decoding instances, we set $L_m$ as the largest batch size.


\parabf{Combat busrtiness.} 
Burstiness in workloads can cause a deluge of KV caches to transfer from prefill to decoding instances, risking memory overload on decoding instances.
To circumvent this, \sys employs a ``pull'' method for KV cache transmission rather than a ``push'' approach -- decoding instances fetch KV cache from prefill instances \emph{as needed}, using the GPU memory of prefill instances as a queuing buffer. This way, the prefill instance can continue handling other prefill jobs by simply retaining the KV Cache in the GPU memory after processing the prompt. Hence, each type of instance operates at its own pace without complex coordination. 


\begin{figure}[!t]
    \centering
    \includegraphics[width=0.9\linewidth]{figures/architecture.pdf}
    \caption{\sys Runtime System Architecture}
    \vspace{-4mm}
    \label{fig:architecture}
\end{figure}

\parabf{Replaning.} The resource and parallelism plan in \sys is optimized for a specific workload pattern, which may become suboptimal if the workload pattern changes over time. \sys implement periodic replanning. A workload profiler monitors key parameters such as the average input and output length of the requests, the average arrival rate, etc. If a significant pattern shift is detected, \sys will trigger a rerun of the placement algorithm based on recent historical data. This process is expedient -- the proposed algorithm runs in seconds (\S\ref{exp:run_time}) and reloading LLM weights can be completed within minutes -- far shorter than the hourly scale at which real-world workload variations tend to occur.



\parabf{Preemption and fault tolerance.} \sys does not implement advanced runtime policies like preemption~\cite{preemption} and fault tolerance~\cite{fault_tolerance}, which are complementary to disaggregation. Nevertheless, we discuss how they fit into \sys.
In \sys, the FCFS policy can lead to a ``convoy effect'', where longer requests block shorter ones in the prefill stage. Incorporating preemptive strategies, as suggested in existing literature \cite{wu2023fast}, could enhance efficiency and is feasible within our system's architecture.
While not a primary focus in the current \sys, fault tolerance is a critical aspect for consideration. In traditional colocation- and replication-based systems, a fault in one instance typically does not disrupt other replica instances. However, in \sys, the dependency between prefill and decoding instances introduces the risk of fault propagation. For example, a fault in a single decoding instance mapped to multiple prefill instances could potentially cripple the entire service and cluster. We leave both as future work.

